\section{Localization and the next quantum induction}
\label{ps19-next}

The homogeneous physical Gaussian is now identified in the source
pivot coordinates, on a common small complex neighborhood with
constants independent of the preceding classical depth. Its
covariance has the spectral and spatial estimates required by
the marked overlap expansion of Notebook 18. The retained
cutoff coordinates are fixed, and the local measure/gauge
factors have their explicit analytic bounds.

The strictly background-local version requires a small local
decomposition of the actual precision and of the action
remainder; exponential matrix-kernel decay alone does not
supply it. Notebook 20 verifies these hypotheses directly for
the bare Wilson step. Notebook 21 constructs them uniformly
through the minimized classical actions using a rooted support
fixed point, and closes one integration of that classical
reference action. Composition with the stationary response
and all inherited quantum corrections must still be included
in a complete induction.

The original CMP116 localization and integration results provide
one route once their orbit/reconstruction and deformed-operator
hypotheses are matched. Notebook 21 supplies the direct
fixed-point route on our homogeneous small product domain.
The next steps must extend to bad regions and running
interactions while preserving the actual normalization.

That induction, uniform sourced observable bounds, and the
continuum symmetry and reconstruction arguments remain the
steps from these finite-regulator results to four-dimensional
Yang--Mills existence. They are the continuing objective.

\begin{thebibliography}{9}
\bibitem{prior19}
\emph{Toward a four-dimensional Yang--Mills construction},
research Notebooks 5--18 in this workspace. The ordinary nonlinear
minimum, classical coarse response and flat retained-coordinate
covariance are the principal inputs.
\bibitem{balaban10919}
T.~Ba\l aban,
\emph{Renormalization group approach to lattice gauge field
theories I}, Commun. Math. Phys. \textbf{109} (1987), 249--301,
\href{https://scispace.com/pdf/renormalization-group-approach-to-lattice-gauge-field-3ygai0vxt1.pdf}{original article copy},
equations (0.11)--(0.16), (2.3)--(2.13).
\bibitem{balaban11619}
T.~Ba\l aban,
\emph{Renormalization group approach to lattice gauge field
theories. II. Cluster expansions},
Commun. Math. Phys. \textbf{116} (1988), 1--22,
\href{https://scispace.com/pdf/renormalization-group-approach-to-lattice-gauge-field-4penm4k7hc.pdf}{original article copy}.
\bibitem{source19}
E.~Sussman and collaborators,
\href{https://github.com/ethanwsussman/Yang--Mills-LLM-experiment}{Yang--Mills-LLM-experiment},
deposited raw accounts at commit
\texttt{e392ef45013a581664bab76b90867170b30fe10e}.
\end{thebibliography}
\end{document}
